%% ===========================================================================
\FloatBarrier
\section{Introduction}
\label{sec:intro}
%% ===========================================================================

The electromechanical response of thin piezoelectric structures depends on
bulk material properties and spatial variations in deformation. Under bending, the strain changes across the thickness, and the
resulting strain gradient contributes to electrical polarization through
flexoelectric coupling. Consequently, the relative piezoelectric and
flexoelectric responses depend on thickness, deformation mode and electrode
configuration. This distinction is relevant to miniature resonators, vibration
energy harvesters and sensing devices~\cite{ErturkInman2011}, as
well as architected materials in which a microstructural length influences the
macroscopic response~\cite{Jiao2023}. Classical piezoelectricity does
not independently represent the strain-gradient contribution. A higher-order
formulation can therefore be used to examine the interaction of strain-gradient
stiffness and electrical coupling in geometries where the characteristic
structural dimensions become comparable to the assumed material length.

The response of a vibrating electromechanical device is also affected by
constitutive dissipation. Frequency-dependent dielectric relaxation has long
been described by distributed-relaxation models~\cite{ColeCole1941,HavriliakNegami1967},
while mechanical hereditary effects can be represented by fractional
viscoelastic constitutive laws. The present study adopts the latter approach:
the fractional derivative acts on a reduced mechanical stiffness operator,
not on the dielectric permittivity or electric displacement. This distinction
is retained in the formulation and in the interpretation of the calculated
ring-downs. Electrical boundary conditions provide another independent source
of variation, since the open- and short-circuit responses depend on the
orientation and location of the electrodes.
Size-dependent mechanical response has been described using
strain-gradient elasticity and nonlocal continuum mechanics. Toupin~\cite{Toupin1962} proposed a theory of
elastic dielectrics in which the energy depends on the second displacement
gradient, Mindlin and Tiersten~\cite{MindlinTiersten1962} formulated a related couple-stress theory, and Mindlin~\cite{Mindlin1964,Mindlin1965} extended
it to a full hierarchy of higher-order elastic constants, first for isotropic
solids and then for crystal classes. Aifantis~\cite{Aifantis1992} proposed to
reduce that hierarchy to a single parameter by keeping only the Laplacian of the
strain in the energy, and Ru and Aifantis~\cite{RuAifantis1993} extended the
reduction to a form that keeps the qualitative predictions of the full theory:
localization is regularized, singularities are removed and stiffness grows as
the member shrinks. In parallel, Eringen and
Edelen~\cite{EringenEdelen1972} replaced the local
stress by an integral average over a neighborhood, so the length entered as the
support of a kernel rather than as a coefficient, and Lim~\cite{Lim2015}
combined the two descriptions in a higher-order nonlocal strain-gradient model
in which an integral kernel and a gradient expansion act at once. In these approaches, the internal length is specified by constitutive
parameters or averaging kernels, with interpretation dependent on the
material microstructure. Discrete-to-continuum relationships provide further context. The review
by Askes and Aifantis~\cite{AskesAifantis2011} compares the formulations, the identification
procedures and the finite element implementations; discrete-to-continuum
arguments show how the gradient terms emerge from a lattice or a representative
volume~\cite{MetrikineAskes2002}; and homogenization of
three-dimensional lattices establishes that a strain-gradient continuum
reproduces size effects that classical homogenization
misses~\cite{Molavitabrizi2023}. The internal-length tensor is often specified as a constitutive parameter,
including in anisotropic gradient models and symmetry classifications
~\cite{Polizzotto2018,Auffray2019}. Connections between integral
averaging and gradient elasticity have also been established by asymptotic
expansion, with attention to higher-order boundary conditions
~\cite{BorinoPolizzotto2014}. Here a particular second-moment tensor is
assigned from a prescribed geometric window and examined inside a coupled
electromechanical discretization; no experimental identification of that
window is claimed.

The higher-order mechanical description is accompanied by flexoelectric
constitutive coupling. Kogan~\cite{Kogan1964} identified the coupling of polarization to the
strain gradient, Mindlin~\cite{Mindlin1968} formulated it as a
polarization-gradient theory, Cross~\cite{Cross2006} reviewed its magnitude in
insulating solids. For barium titanate specifically, Ma and Cross~\cite{MaCross2006}
measured a transverse coefficient $\mu_{12}\simeq\SI{10}{\micro\coulomb\per\meter}$
in unpoled ceramics near room temperature, whereas Hu and co-workers~\cite{HuDengLiangShen2017}
reported a longitudinal coefficient $\mu_{11}\simeq\SI{17.33}{\micro\coulomb\per\meter}$
from shock-wave tests on unpolarized bulk material. These are different
tensor components measured in different specimen and loading conditions;
neither directly identifies the single phenomenological coefficient prescribed
in the present numerical model. Yudin and Tagantsev~\cite{YudinTagantsev2013} and Zubko and
co-workers~\cite{Zubko2013} reviewed its role in thin films. The separate micron-scale bending experiments of Lam and
co-workers~\cite{Lam2003} showed mechanical strain-gradient stiffening;
those experiments did not measure flexoelectric polarization. On the computational side the
formulations divide by how they make the second gradient admissible. Mixed and
staggered methods carry it as an independent field and keep $C^0$
interpolation~\cite{AmanatidouAravas2002,AskesMorataAifantis2008}; micromorphic and virtual element formulations enlarge the
kinematics~\cite{Forest2009,WriggersHudobivnik2023}; Hermite elements provide the required $C^1$ interpolation
~\cite{BognerFoxSchmit1966,PeteraPittman1994,
Papanicolopulos2009}. Coupled piezoelectric--flexoelectric
formulations have already been demonstrated with mixed finite elements,
including two- and three-dimensional models and a recent dynamic plate
formulation~\cite{MaoPurohitAravas2016,SerraoKozinov2024,LiuEtAl2026}.
Related flexoelectric methods include multi-patch isogeometric analysis
~\cite{SharmaAnitescuRabczuk2025}, frequency-domain isogeometric
analysis~\cite{ChenYaoYvonnet2023} and gradient electromechanical finite
elements of arbitrary order~\cite{JoshanSantapuri2024}; benchmarks
with analytical counterparts~\cite{Yang2021} have also been reported. The finite-element choices impose different computational requirements:
mixed formulations introduce additional variables or constraints, whereas
direct $C^1$ elements enforce derivative continuity through their basis
functions.
These are not the only admissible discretizations. Ventura et al.
\cite{VenturaCodonyFernandez2021} enforce derivative continuity weakly with
a $C^0$ interior-penalty method and demonstrate two- and three-dimensional
examples. Codony et al.~\cite{CodonyMarcoFernandezArias2019} use an immersed
hierarchical B-spline formulation with sensing-electrode and corner conditions
on unfitted geometries. These precedents establish coupled flexoelectric
discretization and electrode treatment before the present work. The present
rectangular BFS implementation supplies a different basis for its prescribed
window-moment and mechanical-memory comparisons; no accuracy, cost or
geometric-flexibility advantage over those methods is established here.
Neither approach is uniformly preferable for all geometries. For the strip considered here, rectangular $C^1$ Bogner--Fox--Schmit
elements are combined with reduced-order modal time integration.

Time-dependent constitutive response is described separately.
Fractional viscoelastic models, developed from Caputo-type derivatives
~\cite{Caputo1967} and applied by Bagley and Torvik~\cite{BagleyTorvik1983}
and Koeller~\cite{Koeller1984}, represent power-law mechanical relaxation.
Time-domain damping models~\cite{EnelundLesieutre1999} provide a basis for
representing mechanical dissipation. In related coupled-field problems,
the vibrational response of a size-dependent thermo-piezo-photo-electric
semiconductor has been analyzed under a memory-dependent
Moore--Gibson--Thompson heat-conduction model~\cite{Gupta2024Vibrational}, and a
comparative study~\cite{Gupta2025Piezo} examined models incorporating
piezoelectric, semiconductor and void effects. Those works address thermal and
wave-propagation couplings, rather than validation of the local flexoelectric
strain-gradient finite-element model considered here. The present study instead
examines how fractional \emph{mechanical stiffness} memory affects the
size-dependent electromechanical response. The parameters $\alpha$ and $\tau$ are prescribed rather than fitted to
measured mechanical or dielectric loss. A reduced modal representation is used to evaluate the hereditary convolution
without advancing the complete displacement vector at each time step. The
accuracy of this reduction remains dependent on the retained modes and on the
time discretization. In the condensed electromechanical problem, the
fractional term may act on either the condensed stiffness or the mechanical
block before electrical condensation. These choices define distinct constitutive models, so the matrix
used in the memory term is specified explicitly.

Existing studies cover higher-order electromechanical continua,
nonlocal-to-gradient relations, finite-element methods and fractional
viscoelasticity. This study examines the effects of the prescribed window second moment,
the selected stiffness-memory operator, electrode conditions and strip
thickness on the calculated electromechanical coupling and damping. The contribution is the integrated formulation and its
numerical comparisons for this setting, rather than a claim to a new general
strain-gradient constitutive theory, a new $C^1$ element family or a new
fractional derivative. The numerical analysis addresses the following aspects.

\begin{enumerate}
\item The assumed positive strain-gradient energy is parametrized by the
  in-plane second-moment tensor $M$ of the prescribed window, for which cuboid,
  ellipse and ellipsoid windows give $M_{ii}=a_i^{2}/3$, $a_i^{2}/4$ and
  $a_i^{2}/5$ in closed form. Within the adopted second-moment truncation, two windows
  with the same second moment produce identical solutions: the higher
  moments of the original integral window are not represented. The tensor's
  anisotropy and orientation remain explicit and change the extensional coupling by
  $14$--$17\,\%$.
\item The phenomenological reduced-stiffness memory is a fractional Kelvin--Voigt law of Caputo order
  $\alpha$, discretized by the L1 scheme and convoluted on the reduced modal
  coordinates, with the memory matrix specified explicitly; at $\alpha=1$ the mechanical-only memory gives a damping ratio $0.431$ times as large
  as the condensed one.
\item The coupling factor is
  evaluated from a quasi-static expression that involves no eigenvector, so the
  same measure can be compared across geometries and thicknesses. For the stated
  strip and electrode arrangement, the extensional coupling is $k^{2}=0.5727$ and changes by only $0.07\,\%$ between $10$
  and $1\,\mu$m, whereas the flexoelectric bending coupling grows by a
  factor $24$. Possible parameter-identification measurements and the assumptions of
  the two-dimensional reduction are discussed; the parameters
  are not independently identified from experiments in this study.
\end{enumerate}

The formulation and its limiting cases are followed by the finite-element
method, numerical verification, and a discussion of the computed responses.
Parameter identification and modeling limitations are addressed before the
conclusions.

\FloatBarrier
